\documentclass[11pt]{article}
\usepackage[margin=0.85in]{geometry}
\usepackage[T1]{fontenc}
\usepackage{amsmath,amssymb,booktabs,graphicx,xurl,hyperref}
\graphicspath{{outputs/task1/}{outputs/task2/}{outputs/task3/pusht_train/}{outputs/task3/pusht_train/report_rollouts/media/}{outputs/bonus_flow/pusht_train/}{outputs/bonus_flow/pusht_train/report_rollouts/media/}}

\title{CS 498 HW5 Report\\Diffusion Policy for Push-T}
\author{Your Name (NetID)}
\date{Fall 2026}

\begin{document}
\maketitle

\section{Demonstration sequences and normalization}
Explain how your sequence indices avoid crossing episode boundaries and how
the provided slicing code turns those indices into endpoint padding. Include
\path{outputs/task1/sequence_windows.png} and
\path{outputs/task1/normalization.png}. State the observation/action shapes,
horizon, observation steps, action steps, and normalization range.

\section{DDIM and conditional temporal U-Net}
Write the provided forward-noising equation and your deterministic DDIM
reverse update, including how a skipped previous timestep is handled. Explain
your FiLM modulation and briefly state how the provided embedding and U-Net
use it. Include
\path{outputs/task2/task2_visual_checkpoint.png}; explain how the trajectory
changes as the diffusion timestep increases and report the oracle DDIM reverse
reconstruction error shown in the figure.

\section{Training}
Report the smoke-run outcome and required full run. Include the training and
validation loss curves, model dimensions, batch size, number of gradient steps,
EMA setting, checkpoint interval, training/evaluation devices, and the 50-step DDIM sampler. Explain how normalized action noise and
observation conditioning are used by \path{compute_loss}.

\section{Push-T rollout evaluation}
State the selected checkpoint, selection rule, 20-seed success rate, and mean
maximum reward. Run \path{task3_policy.py --mode visualize} for one successful
seed and the lowest-reward seed, preferring an unsuccessful seed when one
exists. Include both generated \path{rollout_seed_*_frames.png} contact sheets,
identify the two corresponding submitted MP4 filenames, and discuss the
behavior visible in each video. Explain why receding-horizon action chunks can
help or hurt the observed behavior. Confirm that the numerical-backend fields
recorded in \path{eval_metrics.json} match the device used for evaluation.

% Replace the seed values below with your selected seeds.
% \includegraphics[width=\linewidth]{rollout_seed_100000_frames.png}
% \includegraphics[width=\linewidth]{rollout_seed_100007_frames.png}

\section*{Optional bonus: conditional flow matching}
If attempted, state the interpolation path and velocity target, explain the
Euler ODE sampler, report the training configuration and fixed-seed rollout
metrics, and compare its behavior with the required DDIM policy. Include the
bonus loss curve and one generated rollout contact sheet, and name the
corresponding MP4. Delete this section if the bonus was not attempted.

\section*{AI-use declaration}
Check one statement.

\begin{itemize}
  \item[$\square$] I did not use generative AI for this homework.
  \item[$\square$] I used generative AI only in ways explicitly allowed by the
  course policy. Tool/model, purpose, and affected files:
  \underline{\hspace{1.6in}}
\end{itemize}

I confirm that any submitted code, figures, numerical results, and written
analysis comply with the course collaboration and AI-use policy.

\end{document}
